\documentclass{jsarticle}
\usepackage{amsmath}
\usepackage{amssymb}
\begin{document}
\title{慣性の法則と、回転体による反重力 \\
タイミングとともに、別ベクトルとして重力を加える法則}
\author{Masaaki Yamaguchi}
\date{}
\maketitle

   慣性の法則の慣性力が働くときと同時に、
   前もって、タイミングを計り、慣性力とともに、
   逆向きのベクトルとして、慣性力の力と同じエネルギーの
   重力を生成して、慣性力の影響を無くす。この力を反重力という。
   慣性力と逆向きのベクトルとして、反重力として、
   慣性力に加えて、別のベクトルとする。

   反重力の本質は、重力のベクトルの違いであり、
   核エネルギーが特殊相対性理論と一般相対性理論を使って、
   放射性物質吸収体に、この原理を使い、核エネルギーをリサイクルしている。
   これを宇宙人は知っていて、UFOにこの原理を使っている。

   Higgs場は、重力であり、周りを反重力が覆っている。
   一般相対性理論のベクトルの向きを変えた力の装置が、
   ラムダ・ドライバであり、反重力である。

   重力は熱エネルギーである。ベクトルの向きが違う。
   反重力は吸収体である。
   慣性の法則が働いているときに、これが手に入ったが、離れたになり、
   重力をベクトルの向きを変えて、同じエネルギーを加えた場合、
   量子コンピュータが、そのときの差分を使って、反重力が慣性の法則
   を変えて、宙に浮くUFOにしている。 

   まとめると、慣性の法則のときに、重力のベクトルを違う向きとして、
   加えると、手に入ったが、離れたになり、これが反重力である。 

   UFOが現れる前に、ヘリコプターが現れて、そのあとに、UFOが実際に現れる。
   慣性の法則のしばらくして、反重力として、UFOとして、
   宇宙人が人に教えている。

   反重力の生成は、回転体の電磁気力からくるちからであり、ローレンツ力と
   同じでもあり、機体の内部で生成される力であり、そのために、放射性
   物質吸収体を結晶石として、このエネルギーからくる放射力を吸収する
   ために、結晶石に反重力の放射線を吸収する。その上に、外部の
   慣性力を内部の反重力で、手に入ったら、離れたをしている。


\end{document}
